\documentclass[10pt]{article}

%% Verifica: equazioni e disequazioni - classe 3T, scienze applicate (2 ore)
%% (nessun esercizio con il valore assoluto)
%% Opzioni utili:  answers (copia docente con soluzioni), nofigures (bozza veloce)
\usepackage[seed=2026]{fancytests}
\usepackage{lmodern}
\renewcommand*\familydefault{\sfdefault}
\usepackage[T1]{fontenc}

\geometry{a4paper, top=15mm, bottom=15mm, left=17.5mm, right=17.5mm}
\linespread{1.25}

% voto da 1 a 10, punteggio in diciottesimi
\testsetup{class=3T, dsa=none, grade-min=1, grade-max=10, style=compact,
  labels=boxed, mcq=margin}

% indicatori di difficolta': visibili solo nella copia docente (opzione answers)
\makeatletter
\ifft@answers\else\renewcommand*{\ft@diffmarks}[1]{}\fi
% legenda del segno di motivazione: stampata a mano sotto la consegna dell'es. 2
% (la libreria la metterebbe dopo lo spazio di motivazione del primo quesito)
\ExplSyntaxOn
\NewDocumentCommand \legendamotivazione { }
  {
    \bool_gset_true:N \g__ft_legend_done_bool
    \par \nobreak \noindent \hfill { \scriptsize \color { black!55 } \ft@legendtext } \par
  }
\ExplSyntaxOff
\makeatother

% grafici per i quiz
\pgfplotsset{quiz/.style={width=6cm, height=4.3cm, legend style={font=\scriptsize},
  tick label style={font=\scriptsize}}}

\begin{document}
\setseed{203}

\testheader[title={Equazioni e disequazioni}]

%% =========================================================================
\exercise Risolvi le seguenti disequazioni con il metodo che più ritieni opportuno.

\begin{parts}[cols=2, stretch=1.5]
  \itemx[pts=1] $x^3 + 2x^2 - 8x \geq 0$
  \itemx[pts=1] $\sqrt{x + 3} - (x + 1) \leq 0$
  \itemx[pts=1] $(x^4 + 2)(x + 3)(x - 1)^2(2 - x) > 0$
  \itemx[pts=1] $\dfrac{(x + 2)\sqrt{4 - x}}{x^2 - x} \leq 0$
  \itemx[pts=1, diff=3] $\sqrt{x^2 - 2x - 3} > x - 2$
  \itemx[pts=1] $\sqrt{x - 3} + \sqrt{1 - x} \geq 0$
\end{parts}
\exercisescore

\begin{solution}
a) $x(x+4)(x-2) \geq 0$: $-4 \leq x \leq 0 \lor x \geq 2$.
b) $\sqrt{x+3} \leq x + 1$: C.E. $x \geq -3$; serve $x + 1 \geq 0$; elevando al quadrato
$x + 3 \leq x^2 + 2x + 1 \iff x^2 + x - 2 \geq 0 \iff x \leq -2 \lor x \geq 1$. Con $x \geq -1$: $x \geq 1$.
c) $x^4 + 2 > 0$ sempre; $(x-1)^2 > 0$ per $x \neq 1$; conta il segno di $(x+3)(2-x)$:
$-3 < x < 2$, $x \neq 1$.
d) C.E.: $x \leq 4$, $x \neq 0$, $x \neq 1$; $x = 4$ è soluzione (radice nulla). Altrove conta il segno di
$\frac{x+2}{x(x-1)}$: $x \leq -2 \lor 0 < x < 1 \lor x = 4$.
e) Se $x - 2 < 0$ basta la C.E. ($x \leq -1 \lor x \geq 3$): $x \leq -1$. Se $x \geq 2$:
$x^2 - 2x - 3 > x^2 - 4x + 4 \iff x > \frac72$ (compatibile con la C.E.). Soluzione: $x \leq -1 \lor x > \frac72$.
f) Le due radici richiedono $x \geq 3$ e $x \leq 1$ contemporaneamente: nessuna soluzione reale.
\end{solution}

%% =========================================================================
\exercise[motivation=2] Rispondi ai seguenti quesiti (punteggio parziale per risposte errate ma motivate in modo adeguato).
\legendamotivazione

\vspace{.3em}
\begin{parts}
  \itemx[pts=0.5, diff=1] Quale dei grafici rappresenta la funzione $f(x) = \frac{1}{4}(x + 3)(x - 1)(x - 2)^2$?
    \begin{center}
      \begin{ftplot}[quiz, width=4.6cm, height=3.3cm, xmin=-3.9, xmax=3.9, ymin=-4.9, ymax=4.9]
        \fplot[black, very thick, domain=-3.9:3.9, samples=120]{-0.25*(x+3)*(x-1)*(x-2)^2}
        \node[anchor=north west, font=\small\bfseries, fill=white, inner sep=1.5pt] at (rel axis cs:0.03,0.97) {1};
      \end{ftplot}\hfill
      \begin{ftplot}[quiz, width=4.6cm, height=3.3cm, xmin=-3.9, xmax=3.9, ymin=-4.9, ymax=4.9]
        \fplot[black, very thick, dashed, domain=-3.9:3.9, samples=120]{0.25*(x+3)*(x-1)*(x-2)^2}
        \node[anchor=north west, font=\small\bfseries, fill=white, inner sep=1.5pt] at (rel axis cs:0.03,0.97) {2};
      \end{ftplot}\hfill
      \begin{ftplot}[quiz, width=4.6cm, height=3.3cm, xmin=-3.9, xmax=3.9, ymin=-4.9, ymax=4.9]
        \fplot[gray!80, very thick, domain=-3.9:3.9, samples=120]{0.5*(x+3)*(x-1)*(x-2)}
        \node[anchor=north west, font=\small\bfseries, fill=white, inner sep=1.5pt] at (rel axis cs:0.03,0.97) {3};
      \end{ftplot}
    \end{center}
    \begin{choices}[cols=4]
      \correct Il grafico $2$
      \choice Il grafico $1$
      \choice Il grafico $3$
      \choice[fixed] Nessuno dei grafici
    \end{choices}
\end{parts}

\smallskip
\noindent
\begin{minipage}[c]{.35\linewidth}
  \centering
  \begin{tikzpicture}[x=.45cm]
    \draw[-{Stealth[length=2mm]}] (-5.5,0) -- (5.5,0) node[right] {$x$};
    \foreach \t in {-5,...,5} \draw (\t,.08) -- (\t,-.08) node[below] {\scriptsize $\t$};
    \draw[line width=2pt] (-2,0) -- (3,0);
    \fill (-2,0) circle (3pt);
    \draw[thick, fill=white] (3,0) circle (3pt);
  \end{tikzpicture}
\end{minipage}\hfill
\begin{minipage}[c]{.64\linewidth}
  \begin{parts}[ragged]
    \itemx[pts=0.5, diff=2] Quale delle seguenti espressioni ha come insieme di definizione (C.E.) l'intervallo rappresentato in figura?
      \begin{choices}[cols=2]
        \correct $\dfrac{\sqrt{x + 2}}{\sqrt{3 - x}}$
        \choice $\sqrt{x + 2} \cdot \sqrt{3 - x}$
        \choice $\dfrac{\sqrt{x + 2}}{3 - x}$
        \choice $\dfrac{x + 2}{\sqrt{3 - x}}$
      \end{choices}
  \end{parts}
\end{minipage}

\smallskip
\noindent
\begin{minipage}[c]{.35\linewidth}
  \centering
  \begin{ftplot}[quiz, xmin=-4.9, xmax=4.9, ymin=-1.9, ymax=2.4, ytick={-1,0,1,2},
                 minor tick num=1, legend pos=north west]
    \fplot[domain=-4.9:4.9, samples=120]{(x+4)*(x-1)*(x-3)/24}[$g(x)$]
  \end{ftplot}
\end{minipage}\hfill
\begin{minipage}[c]{.64\linewidth}
  \begin{parts}[ragged]
    \itemx[pts=0.5, diff=2] Con $g(x)$ in figura, per quali $x$ si ha $\sqrt{g(x)} \leq 1$?
      \begin{choices}[cols=1]
        \correct $[-4, -3] \cup [-1, 1] \cup [3, 4]$
        \choice $(-4, -3) \cup (-1, 1) \cup (3, 4)$
        \choice $[-4, 1] \cup [3, +\infty)$
        \choice $(-\infty, -3] \cup [-1, 4]$
      \end{choices}
  \end{parts}
\end{minipage}

\smallskip
\noindent
\begin{minipage}[c]{.35\linewidth}
  \centering
  \begin{ftplot}[quiz, xmin=-3.9, xmax=3.9, ymin=-3.9, ymax=4.9, legend style={at={(0.5,0.97)}, anchor=north}]
    \fplot[domain=-3.9:3.9]{x^2-3}[$h(x)$]
    \fplot[domain=-3.9:3.9]{1}[$g(x)$]
    \fplot[domain=-3.9:3.9]{x-1}[$f(x)$]
  \end{ftplot}
\end{minipage}\hfill
\begin{minipage}[c]{.64\linewidth}
  \begin{parts}[ragged]
    \itemx[pts=0.5, diff=2] Utilizzando i grafici in figura, per quali valori di $x$ è verificata la doppia disequazione $f(x) \leq g(x) \leq h(x)$?
      \begin{choices}
        \correct $x \leq -2 \lor x = 2$
        \choice $x \leq -2$
        \choice $-2 \leq x \leq 2$
        \choice $x \leq 2$
      \end{choices}
  \end{parts}
\end{minipage}

\begin{parts}
  \pagebreak
  \itemx[pts=0.5, diff=2] Per quale valore di $k$ l'equazione $x^4 - 4x^2 = k$ ammette esattamente tre soluzioni reali distinte?
    \begin{choices}
      \correct $k = 0$
      \choice $k = -4$
      \choice $k = 4$
      \choice[fixed] Nessun valore di $k$
    \end{choices}
\end{parts}

\smallskip
\noindent
\begin{minipage}[c]{.35\linewidth}
  \centering
  \begin{ftplot}[quiz, xmin=-3.9, xmax=3.9, ymin=-3.9, ymax=3.9, legend pos=north west]
    \fplot[domain=-3.9:3.9, samples=120]{0.25*(x+2)^2*(x-1)}[$f(x)$]
  \end{ftplot}
\end{minipage}\hfill
\begin{minipage}[c]{.64\linewidth}
  \begin{parts}[ragged]
    \itemx[pts=0.5, diff=3] Con $f(x)$ in figura, per quali valori di $x$ si ha $f(x) \cdot (x - 3) \leq 0$?
      \begin{choices}
        \correct $x = -2 \lor 1 \leq x \leq 3$
        \choice $1 \leq x \leq 3$
        \choice $x \leq -2 \lor 1 \leq x \leq 3$
        \choice $-2 \leq x \leq 1 \lor x \geq 3$
      \end{choices}
  \end{parts}
\end{minipage}

\begin{parts}
  \itemx[pts=0.5, diff=3] Per quale $k$ la disequazione $(x - k)^2(x + 2) > 0$ ha come soluzioni $x > -2$, con $x \neq 1$?
    \begin{choices}
      \correct $k = 1$
      \choice $k = -1$
      \choice $k = -2$
      \choice[fixed] Nessun valore di $k$
    \end{choices}
\end{parts}

\smallskip
\noindent
\begin{minipage}[c]{.35\linewidth}
  \centering
  \begin{ftplot}[quiz, xmin=-3.9, xmax=5.9, ymin=-4.9, ymax=4.9, legend pos=north east]
    \fplot[domain=-3.9:5.9]{-0.25*(x+2)*(x-2)*(x-4)}[$g(x)$]
  \end{ftplot}
\end{minipage}\hfill
\begin{minipage}[c]{.64\linewidth}
  \begin{parts}[ragged]
    \itemx[pts=0.5, diff=3] Quale delle seguenti affermazioni sulla funzione $g(x)$ rappresentata in figura è \emph{falsa}?
      \begin{choices}[cols=1]
        \correct $\sqrt{g(x)}$ è definita per $x \in [-2, 2] \cup [4, +\infty)$
        \choice $g(x) \geq 0$ per $x \in (-\infty, -2] \cup [2, 4]$
        \choice $g(x) \cdot (x - 2)^2 \leq 0$ per $x \in [-2, 2] \cup [4, +\infty)$
        \choice $g(0) = -4$
      \end{choices}
  \end{parts}
\end{minipage}

\begin{solution}
a) Zeri semplici in $-3$ e $1$, zero doppio (tangenza) in $2$, coefficiente direttivo positivo.
b) Radice al numeratore: $x \geq -2$; radice al denominatore: $3 - x > 0$.
c) $\sqrt{g} \leq 1 \iff 0 \leq g(x) \leq 1$: $g \geq 0$ su $[-4, 1] \cup [3, +\infty)$, $g \leq 1$ per
$x \leq -3 \lor -1 \leq x \leq 4$.
d) $f \leq g$ per $x \leq 2$; $g \leq h$ per $x \leq -2 \lor x \geq 2$.
e) Con $t = x^2 \geq 0$: $t^2 - 4t - k = 0$; tre soluzioni solo se una radice è $t = 0$ e l'altra positiva, cioè $k = 0$ (soluzioni $0$, $\pm 2$).
f) $f(x)(x-3) = \frac14 (x+2)^2(x-1)(x-3)$: il quadrato non cambia segno ma si annulla in $-2$; conta il segno di $(x-1)(x-3)$.
g) $(x-k)^2 \geq 0$ si annulla solo in $x = k$, che va escluso; serve $k > -2$ e $k = 1$.
h) $\sqrt{g}$ è definita dove $g \geq 0$, cioè su $(-\infty, -2] \cup [2, 4]$.
\end{solution}

%% =========================================================================
\exercise Rispondi ai seguenti quesiti, motivando le risposte.

\begin{parts}[stretch=1.1]
  \itemx[pts=1.5, diff=2] Spiega perché un'equazione polinomiale di quinto grado ammette almeno una e al massimo cinque soluzioni reali. Per ciascuno dei due casi fornisci un esempio di equazione, rappresentando in modo approssimativo il grafico del polinomio corrispondente.
  \itemx[pts=2, diff=3] Rappresenta nello stesso piano il grafico di $f(x) = \sqrt{x + 1}$ e la retta $y = x - 1$ e risolvi graficamente $\sqrt{x + 1} \geq x - 1$, verificando il risultato per via algebrica. Considerando le rette $y = x + k$, determina poi per quali $k$ l'equazione $\sqrt{x + 1} = x + k$ ha esattamente due soluzioni reali distinte.
\end{parts}

\vspace{.4em}
\noindent
\begin{minipage}[c]{.35\linewidth}
  \centering
  \begin{ftplot}[quiz, xmin=-4.9, xmax=4.9, ymin=-4.9, ymax=4.9, legend pos=north east,
                 xtick distance=2, ytick distance=2, minor tick num=1]
    \fplot[domain=-4.9:4.9]{-0.5*(x+3)*(x-1)}[$A(x)$]
  \end{ftplot}
\end{minipage}\hfill
\begin{minipage}[c]{.64\linewidth}
  \begin{parts}[ragged]
    \itemx[pts=1.5, diff=3] In figura è rappresentato il grafico di $A(x)$. Disegna nello stesso piano il grafico di una funzione $B(x)$ tale che la disequazione $\dfrac{A(x)}{B(x)} \leq 0$ abbia come insieme delle soluzioni $(-\infty, -3] \cup [1, 4)$, quindi scrivine l'espressione analitica.
  \end{parts}
\end{minipage}

\begin{solution}
a) Il numero massimo: un polinomio di grado $5$ si scompone in al più $5$ fattori di primo grado.
Il numero minimo: per $x \to +\infty$ e $x \to -\infty$ il polinomio assume segni opposti (decide il
termine $ax^5$), quindi il grafico attraversa almeno una volta l'asse $x$. Esempi: $x^5 + x = x(x^4 + 1) = 0$
(una sola soluzione); $x(x^2 - 1)(x^2 - 4) = 0$ (cinque soluzioni: $0$, $\pm 1$, $\pm 2$).

b) La retta $y = x - 1$ incontra la mezza parabola $y = \sqrt{x+1}$ in $(3, 2)$. Algebricamente: se
$x - 1 < 0$ basta $x \geq -1$, quindi $-1 \leq x < 1$; se $x \geq 1$, $x + 1 \geq x^2 - 2x + 1 \iff
0 \leq x \leq 3$. Soluzione: $-1 \leq x \leq 3$. Rette $y = x + k$: per $k = 1$ la retta passa per
l'origine $(-1, 0)$ della mezza parabola (due soluzioni, $x = -1$ e $x = 0$); la tangenza si ha quando
$x^2 + (2k - 1)x + k^2 - 1 = 0$ ha $\Delta = 5 - 4k = 0$, cioè $k = \frac54$. Due soluzioni per
$1 \leq k < \frac54$.

c) $A(x) > 0$ in $(-3, 1)$, $A(x) < 0$ all'esterno. Serve $B > 0$ per $x < 4$ (così il quoziente è
$\leq 0$ dove $A \leq 0$ e positivo in $(-3, 1)$) e $B < 0$ per $x > 4$, con $B(4) = 0$ (per escludere
$x = 4$). Ad es. $B(x) = 4 - x$.
\end{solution}

%% =========================================================================
\exercise Una semicirconferenza ha diametro $AB = \qty{10}{cm}$. Un punto $P$ si muove sulla
semicirconferenza e $H$ è la sua proiezione sul diametro; poni $AH = x$ (in cm). Esprimi le lunghezze
delle corde $PA$ e $PB$ in funzione di $x$ e determina per quali posizioni di $P$ il perimetro del
triangolo $APB$ è di almeno $\qty{24}{cm}$.
\exercisescore[3]

\begin{solution}
$0 \leq x \leq 10$. Il triangolo $APB$ è rettangolo in $P$ (inscritto in una semicirconferenza) e
$PH = \sqrt{AH \cdot HB} = \sqrt{x(10 - x)}$; per Pitagora nel triangolo $AHP$:
$PA = \sqrt{x^2 + 10x - x^2} = \sqrt{10x}$ e analogamente $PB = \sqrt{10(10 - x)} = \sqrt{100 - 10x}$.
Perimetro: $\sqrt{10x} + \sqrt{100 - 10x} + 10 \geq 24 \iff \sqrt{10x} + \sqrt{100 - 10x} \geq 14$.
Entrambi i membri sono positivi: $100 + 2\sqrt{10x(100 - 10x)} \geq 196 \iff \sqrt{10x(100 - 10x)} \geq 48
\iff 100x^2 - 1000x + 2304 \leq 0 \iff x^2 - 10x + 23.04 \leq 0 \iff 3.6 \leq x \leq 6.4$.
Agli estremi le corde misurano $6$ e $8$ (triangoli $6$-$8$-$10$).
\end{solution}

%% =========================================================================
\bonusexercise Per quali valori del parametro $k \in \spR$ la disequazione $\sqrt{x^2 + 1} - kx > 0$ è verificata per ogni $x \in \spR$?
\bonusexercisescore[1]
\vspace{-.2em}
\bonusexercisenote

\begin{solution}
Per $-1 \leq k \leq 1$: $kx \leq \abs{k}\abs{x} \leq \abs{x} < \sqrt{x^2 + 1}$, quindi la disequazione
vale per ogni $x$. Se $k > 1$, per $x > 0$: $\sqrt{x^2+1} < kx \iff x^2 + 1 < k^2x^2 \iff
x > \frac{1}{\sqrt{k^2 - 1}}$, quindi esistono $x$ per cui è falsa (analogamente per $k < -1$, con
$x < 0$). Graficamente: le rette $y = kx$ restano sotto il grafico di $\sqrt{x^2+1}$ solo se la loro
pendenza è compresa tra $-1$ e $1$. Soluzione: $-1 \leq k \leq 1$.
\end{solution}

\vfill
\finalscore
\end{document}
